\section{Graphs of Diameter Two}
\label{sec:diameter-two}

The first class in which the problem can be solved efficiently consists of the graphs of diameter two, a class studied in algebraic graph theory~\cite{GodsilRoyle2001}.

\begin{definition}[Diameter]
\label{def:diameter}
A graph $G = (V, E)$ has \emph{diameter at most $k$} if $d(x, y) \leq k$ for every pair of vertices $x, y \in V$.
The smallest such $k$ is the \emph{diameter} of $G$, denoted by $\diam(G)$.
\end{definition}

We derive a formula for $\cp(u,v)$ in graphs of diameter two; it yields a polynomial-time algorithm (Remark~\ref{rem:algorithm}).
The proof uses a decomposition: in any graph, a cut $C$ induces a partition of the vertex set into four subsets $I, J, K, L$.
Each vertex lies on one of the two sides $A_1, A_2$ of the cut, and it is either incident to an edge of the cut or not.

\begin{lemma}
\label{lem:cut-decomposition}
Let $G = (V, E)$ be a graph, and let $C$ be a cut that divides $G$ into two sets $A_1$ and $A_2$ (i.e., $V = A_1 \cup A_2$, $A_1 \cap A_2 = \emptyset$, and $C$ is the set of all edges between $A_1$ and $A_2$).
Then there exists a unique decomposition of $A_1$ and $A_2$ into subsets $I, J, K, L$ such that:
\begin{itemize}
    \item $I \cup J = A_1$ and $I \cap J = \emptyset$,
    \item $K \cup L = A_2$ and $K \cap L = \emptyset$,
    \item $\forall x \in J \cup K, \, \exists y \in V: \{x, y\} \in C,$
    \item $\forall x \in I \cup L, \, \nexists y \in V: \{x, y\} \in C.$
\end{itemize}
\end{lemma}

In a graph of diameter two, the vertices that are not incident to the cut lie on at most one side of it.

\begin{lemma}
\label{lem:empty-i-or-l}
Let $G = (V, E)$ be a graph of diameter two, let $C$ be a cut that divides $G$ into two sets $A_1$ and $A_2$, and let $I, J, K, L$ be the decomposition of $A_1$ and $A_2$ described in Lemma~\ref{lem:cut-decomposition}.
Then at least one of the sets $I$ and $L$ is empty:
\[
I = \emptyset \quad\text{or}\quad L = \emptyset.
\]
\end{lemma}

\begin{proof}
Suppose, for a contradiction, that there are vertices $x \in I$ and $y \in L$.
Since $x \in A_1$, $y \in A_2$ and $C$ is the set of all edges between $A_1$ and $A_2$, every $x$--$y$ path contains an edge $\{a, b\}$ of $C$ with $a \in A_1$ and $b \in A_2$.
The vertices $a$ and $b$ are incident to an edge of $C$, so $a \in J$ and $b \in K$; in particular, $a \neq x$ and $b \neq y$, because $x$ and $y$ are not incident to any edge of $C$.
Hence the path has at least one edge between $x$ and $a$, the edge $\{a, b\}$, and at least one edge between $b$ and $y$; that is, it passes through $J$ and $K$ and has at least three edges.
Therefore $d(x, y) \geq 3$, which contradicts the assumption that $G$ has diameter two.
\end{proof}

The next lemma holds in every graph; it concerns the intersection of a cut and a path.

\begin{lemma}
\label{lem:odd-intersection}
Let $G = (V, E)$ be a graph, and let $u, v \in V$ be two distinct vertices.
Let $C$ be a cut that divides $G$ into two sets $A_1$ and $A_2$ with $u \in A_1$ and $v \in A_2$, and let $P$ be a path that connects $u$ and $v$.
Then the path $P$ contains an odd number of edges of $C$:
\[
|C \cap P| \equiv 1 \pmod{2}.
\]
\end{lemma}

\begin{proof}
Let $S = C \cap P$ be the set of edges in the intersection of the cut and the path.
When we traverse the path $P$ from $u$ to $v$, each edge of $S$ changes the side of the cut (i.e., it switches between $A_1$ and $A_2$), while every other edge of $P$ stays on the same side.
The path starts in $A_1$ and ends in $A_2$.
If the path crossed the cut an even number of times, its first and its last vertex would lie on the same side, which contradicts the assumption that $u \in A_1$ and $v \in A_2$.
Therefore, the path crosses the cut an odd number of times, i.e., $|S|$ is odd.
\end{proof}

Lemmas~\ref{lem:empty-i-or-l} and~\ref{lem:odd-intersection} give the value of $\cp(u,v)$ in graphs of diameter two.

\begin{theorem}
\label{thm:diameter-two}
Let $G = (V, E)$ be a connected graph of diameter two, and let $u, v \in V$ be two distinct vertices.
Then a minimum cut-path between $u$ and $v$ has
\[
\cp(u, v) = c(u, v) + d(u, v) - 1
\]
edges.
\end{theorem}

\begin{proof}
We distinguish two cases based on the values of $d(u, v)$ and $c(u, v)$.
Since $G$ has diameter two, $d(u, v) \in \{1, 2\}$; since $G$ is connected, $c(u, v) \geq 1$.
Hence the two cases below cover all pairs $u, v$.

\emph{Case 1: $d(u, v) = 1$ or $c(u, v) = 1$.}
By Lemma~\ref{lem:basic-bounds}, in this case we have
\[
\cp(u, v) = c(u, v) + d(u, v) - 1.
\]

\emph{Case 2: $d(u, v) = 2$ and $c(u, v) \geq 2$.}
By Lemma~\ref{lem:basic-bounds},
\[
c(u, v) \leq \cp(u, v) \leq c(u, v) + 2 - 1 = c(u, v) + 1.
\]
Suppose, for a contradiction, that $\cp(u, v) = c(u, v)$; by the bounds above, this is the only alternative to $\cp(u, v) = c(u, v) + 1$.
Then a minimum cut-path $S$ has the same size as a minimum cut.
The set $S$ contains a $u$--$v$ cut, which has at least $c(u, v) = |S|$ edges, so this cut is all of $S$.
Hence the set $S$ itself is a minimum cut, which we denote by $C$, and the $u$--$v$ path $P$ contained in $S$ is a subset of this cut.

The cut $C$ is the set of all edges between two sides with $u$ and $v$ on different sides.
Indeed, let $A$ be the vertex set of the component of $G \setminus C$ that contains $u$.
Every edge between $A$ and $V \setminus A$ belongs to $C$, since otherwise its end outside $A$ would lie in the same component as $u$.
These edges separate $u$ from $v$, because $v \notin A$; by the minimality of $C$, there are at least $|C|$ of them, so they form all of $C$.
We denote the sides $A$ and $V \setminus A$ by $A_1$ and $A_2$.

By Lemma~\ref{lem:odd-intersection}, applied to $C$ and to the path $P \subseteq C$, the number of edges of $P$ is odd.
Since $d(u,v) = 2$, the path $P$ has at least two edges, and hence at least three.
Let $I, J, K, L$ be the decomposition induced by the cut $C$ (Lemma~\ref{lem:cut-decomposition}).
Two exchanges do not affect the argument.
The first exchanges the names $A_1$ and $A_2$, and with them $I$ with $L$ and $J$ with $K$.
The second exchanges $u$ and $v$: the values $c(u, v)$, $d(u, v)$ and $\cp(u, v)$ do not depend on the order of $u$ and $v$, and $P$ read backwards is a $v$--$u$ path.
By Lemma~\ref{lem:empty-i-or-l} and the first exchange, we may assume that $L = \emptyset$; by the second exchange, we may then assume that $u \in A_2$.
Since $A_2 = K \cup L = K$, we have $u \in K$, so the graph is partitioned into the three sets $I, J, K$.
This agrees with the fact that the first edge of $P$ belongs to the cut and is incident to $u$.
We observe that $v \in J$: the last edge of $P$ belongs to the cut and is incident to $v$, so the vertex $v$ cannot belong to the set $I$.
This structure is illustrated in Figure~\ref{fig:diameter-two-structure}.

\begin{figure}[H]
    \centering
    \includegraphics{diameter-two-structure.pdf}
    \caption{Structure of the cut $C$ in the proof of Theorem~\ref{thm:diameter-two}; the edges counted in the estimate are drawn in red}
    \label{fig:diameter-two-structure}
\end{figure}

Every edge of $C$ has exactly one end in $K$, because $C$ is the set of all edges between $A_1 = I \cup J$ and $A_2 = K$.
Counting the edges of $C$ according to their end-vertex in $K$, we estimate the size of the cut $C$:
\[
\cp(u, v) = c(u,v) = |C| \geq \deg(u, I) + \deg(u, J) + 2 + (|K| - 2).
\]
Here $\deg(u, I) = 0$, because $u \in K$ has no neighbors in $I$: an edge between $u \in A_2$ and a vertex of $I \subseteq A_1$ would belong to $C$, but no vertex of $I$ is incident to an edge of $C$.
The terms of the estimate count the following edges:
\begin{itemize}
    \item $\deg(u, I)$ and $\deg(u, J)$ denote the numbers of edges that lead from the vertex $u$ to the opposite side of the minimum cut,

    \item the path $P$ is contained in the cut and has at least three edges, so its vertices alternate between $A_2$ and $A_1$; hence its third vertex lies in $K$, differs from $u$ because a path does not repeat vertices, and contributes at least two edges to the cut, namely the second and the third edge of $P$,

    \item each of the remaining $|K| - 2$ vertices of the set $K$ is incident to an edge of $C$ by the definition of $K$, so it contributes at least one edge to the cut.
\end{itemize}
The edges counted on the right-hand side are drawn in red in Figure~\ref{fig:diameter-two-structure}.

Since $\deg(u, I) \geq 0$, and since $\deg(u, K) \leq |K| - 1$ because $u \in K$ is not its own neighbor, we obtain
\[
\deg(u, I) + \deg(u, J) + 2 + (|K| - 2) \geq
0 + \deg(u, J) + \deg(u, K) + 1 = \deg(u) + 1 \geq c(u, v) + 1.
\]
The equality holds because $V = I \cup J \cup K$ and $u$ has no neighbors in $I$.
The last inequality holds because the edges incident to $u$ form a $u$--$v$ cut.
This contradicts the assumption that $\cp(u, v) = c(u, v)$; hence $\cp(u, v) = c(u, v) + 1 = c(u, v) + d(u, v) - 1$.
\end{proof}

\begin{remark}
\label{rem:algorithm}
Whenever $\cp(u, v) = c(u, v) + d(u, v) - 1$, the proof of Lemma~\ref{lem:basic-bounds} shows that the union of an arbitrary minimum $u$--$v$ cut and an arbitrary shortest $u$--$v$ path is a minimum cut-path.
Hence, in graphs of diameter two, a minimum cut-path can be found by one minimum-cut computation and one breadth-first search.
\end{remark}

\section{Graphs with Cut-Value at Most Two}
\label{sec:cut-two}

The same formula holds in graphs in which the minimum cut-value of every pair of vertices is at most two.

Such graphs have the structure of a \emph{cactus graph} (cf.~\cite{Mehlhorn2017Certifying}): they consist of cycles and single edges, any two of which share at most one vertex, an articulation point.
Indeed, suppose that two distinct cycles shared an edge.
Then the second cycle contains a path whose ends $x \neq y$ lie on the first cycle and whose edges do not belong to the first cycle.
The union of the two cycles would then contain the two vertices $x$ and $y$ joined by three edge-disjoint paths: this path and the two $x$--$y$ paths along the first cycle.
Separating $x$ from $y$ would require at least three edges, contradicting the assumption $c(x,y) \leq 2$.

\begin{theorem}
\label{thm:cut-two}
Let $G = (V, E)$ be a connected graph such that $c(x, y) \leq 2$ for every pair of distinct vertices $x, y \in V$.
Then, for any two distinct vertices $u,v \in V$, a minimum cut-path between $u$ and $v$ has
\[
\cp(u, v) = c(u, v) + d(u, v) - 1
\]
edges.
\end{theorem}

\begin{proof}
We distinguish two cases based on the values of $d(u, v)$ and $c(u, v)$.
Since $G$ is connected, $c(u, v) \geq 1$, and $c(u, v) \leq 2$ by assumption; since $u \neq v$, $d(u, v) \geq 1$.
Hence the two cases below cover all pairs $u, v$.

\emph{Case 1: $d(u, v) = 1$ or $c(u, v) = 1$.}
The formula holds by Lemma~\ref{lem:basic-bounds}.

\emph{Case 2: $d(u, v) \geq 2$ and $c(u, v) = 2$.}
By Lemma~\ref{lem:basic-bounds}, we have
\[
d(u, v) \leq \cp(u, v) \leq d(u, v) + 1.
\]
Suppose, for a contradiction, that $\cp(u, v) = d(u, v)$.
A minimum cut-path $S$ then has $d(u, v)$ edges and contains a $u$--$v$ path, which has at least $d(u, v)$ edges, so $S$ is this path.
Then a minimum cut-path is a shortest $u$--$v$ path $P$ that contains a $u$--$v$ cut, so the graph $G \setminus P$ contains no $u$--$v$ path.

Since $c(u, v) = 2$, no single edge separates $u$ from $v$.
A bridge of $G$ on the path $P$ would separate $u$ from $v$; hence no edge of $P$ is a bridge, and every edge of $P$ lies on a cycle of the cactus graph $G$.
The path $P$ therefore passes through a sequence of cycles $Z_1, Z_2, \dots, Z_t$, and in each cycle $Z_i$ it follows one of the two arcs between the vertex where it enters $Z_i$ and the vertex where it leaves $Z_i$.
The path $P$ enters $Z_1$ at $u$ and leaves $Z_t$ at $v$.
For $i < t$, the vertex where $P$ leaves $Z_i$ is the vertex where it enters $Z_{i+1}$, since every edge of $P$ lies on one of these cycles.
The other arc of each cycle $Z_i$ is edge-disjoint from $P$, and consecutive arcs share these vertices, so the concatenation of these arcs connects $u$ and $v$ in $G \setminus P$.
This contradicts the fact that $G \setminus P$ contains no $u$--$v$ path.
Therefore,
\[
\cp(u, v) = d(u, v) + 1 = c(u, v) + d(u, v) - 1. \qedhere
\]
\end{proof}

By Remark~\ref{rem:algorithm}, a minimum cut-path in such a graph is again obtained as the union of a minimum cut and a shortest path.
